%\documentclass{book}\usepackage{sjtfonts}\usepackage{includegraphics}\usepackage{alltt}\usepackage{latexsym}
%\usepackage{array}\usepackage{multicol}
%\begin{document}\input{defs0}%		miradefs.tex 		    June 1994

% opening and closing a quoted Miranda section

\newcommand{\mira}{\noindent\hrulefill\nopagebreak\so}
\newcommand{\arim}{\nopagebreak\st\hrulefill\\}

% backslash in the Miranda environment
% Only feasible approach is to use \verb+\+
% in line!

%\newcommand{\bs}{$\mi\backslash\!\!$}
%\newcommand{\lbs}{\mbox{\verb+\+}}

% ignoring part of the input

\newcommand{\ignore}[1]{}

% a blank line in a Miranda section

\newcommand{\bl}{\mbox{\ }}

% Signalling a diffculty.

\definecolor{light-gray}{gray}{0.95}

\newcommand{\beware}[2]
 {\medskip\noindent\colorbox{light-gray}{\parbox{\textwidth}
 {\mbox{\large\textbf{#1}} \medskip\noindent\\ #2}\medskip\noindent}}

% Subscripting in Miranda text
\newcommand{\subscr}[2]{\mbox{${\mbox{\texttt{#1}}}_{\mbox{\texttt{#2}}}$}}

%\newcommand{\subscr}[2]{\mbox{${\mbox{\mi #1}}_{\mbox{\smtt #2}}$}}
\newcommand{\aone}{\subscr{a}{1}}
\newcommand{\atwo}{\subscr{a}{2}}
\newcommand{\ak}{\subscr{a}{k}}

\newcommand{\aoneone}{\subscr{a}{1,1}}

\newcommand{\eone}{\subscr{e}{1}}
\newcommand{\etwo}{\subscr{e}{2}}
\newcommand{\ei}{\subscr{e}{i}}
\newcommand{\er}{\subscr{e}{r}}

\newcommand{\fone}{\subscr{f}{1}}
\newcommand{\fl}{\subscr{f}{l}}

\newcommand{\gone}{\subscr{g}{1}}
\newcommand{\gtwo}{\subscr{g}{2}}
\newcommand{\gi}{\subscr{g}{i}}
\newcommand{\gr}{\subscr{g}{r}}

\newcommand{\hone}{\subscr{h}{1}}
\newcommand{\hl}{\subscr{h}{l}}

\newcommand{\pone}{\subscr{p}{1}}
\newcommand{\ptwo}{\subscr{p}{2}}
\newcommand{\pk}{\subscr{p}{k}}

\newcommand{\qone}{\subscr{q}{1}}
\newcommand{\qtwo}{\subscr{q}{2}}
\newcommand{\qk}{\subscr{q}{k}}

\newcommand{\rone}{\subscr{r}{1}}
\newcommand{\rtwo}{\subscr{r}{2}}

\newcommand{\tone}{\subscr{t}{1}}
\newcommand{\ttwo}{\subscr{t}{2}}
\newcommand{\tn}{\subscr{t}{n}}

\newcommand{\vone}{\subscr{v}{1}}
\newcommand{\vtwo}{\subscr{v}{2}}
\newcommand{\vn}{\subscr{v}{n}}

% for regular expressions

\newcommand{\stone}{\subscr{st}{1}}
\newcommand{\sttwo}{\subscr{st}{2}}
\newcommand{\stn}{\subscr{st}{n}}
\newcommand{\rthree}{\subscr{r}{3}}

\newcommand{\eps}{$\varepsilon$}
\newcommand{\trans}[1]{\stackrel{\mbox{\tt #1}}{\longrightarrow}}



% superscript in Miranda text

\newcommand{\superscr}[2]{\mbox{${\mbox{\mi #1}}^{\mbox{\smtt #2}}$}}

% Not equals in Miranda text

\newcommand{\noteq}{\makebox[0pt][l]{=}/}
\newcommand{\overslash}{\makebox[0pt][l]{/}}

% Definitions in complexity chapter

\newfont{\oldSym}{cmsy10}
\newcommand{\Th}{\boldmath$\oldSym\Theta$}
\newcommand{\pow}[1]{\^{}#1}
\newcommand{\leqq}{$\leq$}
\newcommand{\geqq}{$\geq$}
\newcommand{\lll}{$\ll$}
\newcommand{\eqq}{$\equiv$}
\newcommand{\ltwo}{\subscr{log}{2}}

% Indexing

\newcommand{\minx}[1]{\index{#1@{\ttfamily{}#1}}}
\appendix


%\minitocoptionfalse
\chapter{GHCi errors}
\label{hugsErrors}
\index{errors|(}\index{GHCi!errors|(}

This appendix examines some of the more common programming errors in Haskell,
and shows the error messages to which they give rise in GHCi.

The programs we write all too often contain errors. On encountering an error,
the system either halts, and gives an \textbf{error message\/}, or
continues, but gives a \textbf{warning message\/} to tell us that something
unusual has happened, which might signal that
we have made an error. In this
appendix, we look at a selection of the messages output by GHCi;
we have
chosen the messages which are both common and require some explanation;
messages like
\begin{alltt}
*** Exception: Prelude.head: empty list
\end{alltt}
are self-explanatory. The messages are classified into roughly distinct
areas. Syntax errors show up malformed programs, while type errors show
well-formed programs in which objects are used at the wrong types. In fact,
an ill-formed expression can often show itself as a type error and not as a
syntax error, so the boundaries are not clear.

\section*{Syntax errors}
\index{errors!syntax}\index{syntax error}

A Haskell system attempts to match the input we give to the
syntax of the language. Commonly, when something goes wrong, we type
something \emph{unexpected}. 
\begin{itemize}

\item
Typing `\texttt{2==3)}' will provoke the error
message
\begin{alltt}
<interactive>:1:4: parse error on input `)'
\end{alltt}
\item
If a part of a definition is missing, as in
\begin{alltt}
fun x
fun 2 = 34
\end{alltt}
we receive the message
\begin{alltt}
Errors.hs:3:0: Not in scope: `fun'
Errors.hs:3:4: Not in scope: `x'
\end{alltt}
The problem here is that the system tries to understand \texttt{fun} and \texttt{x} and these are not (yet) defined. 

\item
A similar error results when the two lines are reversed, as in 
\begin{alltt}
fun 2 = 34
fun x
\end{alltt}
when the error message is 
\begin{alltt}
Errors.hs:4:4: Not in scope: `x'
\end{alltt}

\item
The inclusion of a type definition in a \texttt{where}
clause, like this
\begin{alltt}
fun x = x+1
        where
          type MyInt = Int
\end{alltt}
is signalled by 
\begin{alltt}
Errors.hs:8:10: parse error on input `type'
\end{alltt}

\item
The syntax of patterns is more restricted
than the full expression syntax, and so we get error messages like
\begin{alltt}
Errors.hs:6:5:
    Conflicting definitions for `x'
    Bound at: Errors.hs:6:5
              Errors.hs:6:7
    In the definition of `fun'
\end{alltt}
when we use the same variable more than once within a pattern, as in the definition
\begin{alltt}
fun (x,x) = x+1
\end{alltt}

\item
In specifying constants, we can make errors: floating-point numbers can be
too large, and characters specified by an out-of-range ASCII code: for example, on typing \texttt{'\bs{}9999999999999999'} as input we get
\begin{alltt}
<interactive>:1:18:
    lexical error in string/character literal at character '\bs{}''
    \end{alltt}

\item
Not every string can be used as a name; some words in Haskell are
\textbf{keywords} or \textbf{reserved identifiers},
and will give an error if used as an identifier. The keywords
are
\index{keyword}
\index{identifier!reserved}
\begin{alltt}
case class data default deriving do else if import in infix
infixl infixr instance let module newtype of then type where
\end{alltt}
For example, including the definition 
\begin{alltt}
do (x,y) = x+1
\end{alltt}
gives this error message
\begin{alltt}
Errors.hs:6:0: Parse error in pattern
\end{alltt}
and the definition
\begin{alltt}
data (x,y) = x+1
\end{alltt}
gives this message
\begin{alltt}
Errors.hs:6:13: Not a data constructor: `x'
\end{alltt}
As you can see from the two examples, the error message in a case like this reflects the  meaning of the keyword.

\item
The special identifiers \texttt{as}, \texttt{qualified} and
\texttt{hiding} have special meanings in certain contexts but can be
used as ordinary identifiers.

\item
The final restriction on names is that names of constructors and types
must begin with a
capital letter; nothing else can do so, and hence we get error messages like
\begin{alltt}
Errors.hs:6:0: Not in scope: data constructor `Montana'
\end{alltt}
if we try to define a function called \texttt{Montana}.
\end{itemize}

\section*{Type errors}
\index{errors!type}\index{type error}

In this section we look at various different type errors that we can provoke in GHCi.

\begin{itemize}
\item  
As we have seen in the body of the text, the main type error we meet is
exemplified by the response to typing \texttt{'c' \&\& True} to the GHCi prompt:
\begin{alltt}
Couldn't match expected type `Bool' against inferred type `Char'
In the first argument of `(\&\&)', namely 'c'
In the expression: 'c' \&\& True
In the definition of `it': it = 'c' \&\& True
\end{alltt}
which is provoked by using a \texttt{Char} where an \texttt{Bool} is expected. 
\item
Other type errors, such as
\begin{alltt}
True + 4        
\end{alltt}
provoke the error message
\begin{alltt}
No instance for (Num Bool)
  arising from a use of `+' at <interactive>:1:0-7
Possible fix: add an instance declaration for (Num Bool)
In the expression: True + 4
In the definition of `it': it = True + 4
\end{alltt}
This comes from the class mechanism: the system attempts to make \texttt{Bool}
an instance of the class \texttt{Num} of numeric types over which `\texttt{+}' is
defined. The error results since there is no such instance
declaration making \texttt{Bool} belong to the class \texttt{Num}.
\item 
As we said before,
we can get type errors from syntax errors. For example, writing \texttt{abs -2}
instead of \texttt{abs (-2)} gives the error message
\begin{alltt}
No instance for (Num (a -> a))
  arising from a use of `-' at <interactive>:1:0-5
Possible fix: add an instance declaration for (Num (a -> a))
In the expression: abs - 2
In the definition of `it': it = abs - 2
\end{alltt}
because it is parsed as \texttt{2} subtracted from \texttt{abs::a->a}, and
the operator `\texttt{-}' expects something in the class \texttt{Num}, 
rather than a function of type
\texttt{a->a}. Other common type errors come from confusing the roles of `\texttt{:}'
and `\texttt{++}' as in \texttt{2++[2]} and \texttt{[2]:[2]}.
\item
We always give type declarations for our definitions; one advantage of this is
to spot when our definition does not conform to its declared type. For
example,
\begin{alltt}
myCheck :: Char -> Bool
myCheck n = toEnum n == 6
\end{alltt}
gives the error message
\begin{alltt}
Couldn't match expected type `Int' against inferred type `Char'
In the first argument of `toEnum', namely `n'
In the first argument of `(==)', namely `toEnum n'
In the expression: toEnum n == 6
\end{alltt}
Without the type declaration the definition would be accepted, only to give
an error (presumably) when it is used.
\item
A definition like
\begin{alltt}
asc x y 
  | x <= y  = x y
\end{alltt}
will give this error:
\begin{alltt}
Occurs check: cannot construct the infinite type: a = a -> t
Probable cause: `x' is applied to too many arguments
In the expression: x y
In the definition of `asc': asc x y | x <= y = x y
\end{alltt}
The problem here is that \texttt{x} is \emph{compared with} \texttt{y} in the guard (type \texttt{a}), but \emph{applied to} \texttt{y} in the body (type \texttt{a->t}): the unification required here produces an infinite type, which is not allowed.
\item
A final error related to types is given by definitions like
\begin{alltt}
type Fred = (Fred,Int) \hfill (Fred)
\end{alltt}
a \textbf{recursive\/} type synonym; these are signalled by
\begin{alltt}
Cycle in type synonym declarations:
  Errors.hs:9:0-21: type Fred = (Fred, Int)
\end{alltt}
The effect of \texttt{(Fred)} can be modelled by the algebraic type definition
\begin{alltt}
data Fred = Node Fred Int
\end{alltt}
which introduces the \textbf{constructor\/} \texttt{Node} to identify objects of
this type.
\end{itemize}


\section*{Program errors}
\index{errors!program}\index{program error}

Once we have written a syntactically and type
correct script, and asked for the value of an
expression which is itself acceptable, other errors can be produced during
the \textbf{evaluation\/} of the expression.
\begin{itemize}
\item
The first class of errors comes from missing cases in definitions. If we
have written a definition like
\begin{alltt}
bat [] = 45
\end{alltt} 
and applied it to \texttt{[34]} we get the response
\begin{alltt}
*** Exception: Errors.hs:11:0-10: Non-exhaustive patterns 
               in function bat
\end{alltt}
which shows the point at which evaluation can go no further, since there is no
case in the definition of \texttt{bat} to cover a non-empty list. Similar errors
come from built-in functions, such as \texttt{head}.

\item
Other errors happen because an \textbf{arithmetical constraint\/} has been
broken. These include an out-of-range list index, division by zero, using a
fraction where an integer is expected and floating-point calculations which go
out of range; the error messages all have the same form. For example, suppose that we evaluate 
\begin{alltt}
3 `div` 0
\end{alltt}
then the error is
\begin{alltt} 
*** Exception: divide by zero
\end{alltt}

\item
If we make a conformal definition, like
\begin{alltt}
[a,b] = [1 ..\ 10]
\end{alltt}
this will fail with the message
\begin{alltt}
*** Exception: Errors.hs:13:0-16: Irrefutable pattern failed 
               for pattern [a, b]
\end{alltt}
\noindent
when either \texttt{a} or \texttt{b} is evaluated.

\end{itemize}
\vspace{-8pt}

\section*{Module errors}
\index{errors!module}\index{modules!errors}
\vspace{-5pt}

The \texttt{module} and \texttt{import} statements can provoke a variety of
error messages: files may not be present, or may contain errors; names may be
included more than once, or an alias on inclusion may cause a name clash. The
error messages for these and other errors are self-explanatory.
\vspace{-8pt}

\section*{System messages}
\index{system messages}
\vspace{-5pt}

In response to some commands and interrupts, the system generates messages,
including
\begin{itemize}
\item
\begin{alltt}
\^{}C ... Interrupted
\end{alltt}
signalling the interruption of the current task by typing \texttt{Ctrl-C}.
\item
The message
\begin{alltt}
<interactive>: memory allocation failed (requested 2097152 bytes)
\end{alltt}
which shows that the space consumption of the evaluation exceeds that
available. One way around this is to increase the size of the heap. 
\end{itemize}

A measure of the space complexity of a function, as described in Chapter
\ref{behaviour}, is given by the size of the smallest heap in which the
evaluation can take place; how this `residency' is measured is described in that chapter.\index{errors|)}\index{GHCi!errors|)}

%\end{document}
